\scanpage{17}
\uline{Remarque} Supposons que $F$, en plus de a) b) c) \deleted{(et (Dim))}\inserted{d) e)} satisfasse encore f). Alors il existe une décomposition cellulaire \uline{stricte} \inserted{$\mathcal P$} (i.e.: feuilles connexes) $F$-admissible telle que les $Y\in\Phi$ soient des \inserted{réunions} feuilles, et \unsure{$=$|---} parmi ces $\mathcal P$, il y en a une plus grossière. La dém. est la \unsure{même}, mais au lieu de prendre les $U_i$ précédents, on prend leurs composantes connexes.

Donc $F$ est réunion filtrante de $F_{\mathcal P}$ avec les $\mathcal P$ décompositions cellulaires à \uline{feuilles connexes}; et ces décompositions cellulaires admettent des sup. finis.

\uline{§ 4} \inserted{(Décompositions cellulaires $F$-modérées \uline{régulières})}
Supposons maintenant donné une partie $F_r$ de $F$ \deleted{\unsure{les|elle}} formée d'éléments de $F$ \deleted{\unsure{tous}}\inserted{loc.} fermés: ce seront les parties $F$-\uline{régulières} de $F$. On suppose que

\textcircled{$\unsure{B|S}$} Pour tout $Y\in F$, $Y$ \deleted{\unsure{loc.}} fermé, parmi les parties \inserted{rel.} ouvertes $U$ de $Y$ qui sont $\in F_r$, il y en a une plus grande $Y_{\mathrm{reg}}$, et \deleted{\unsure{Y}} celle-ci est dense dans $Y$ i.e. son complémentaire $Y-Y_{\mathrm{reg}}=Y_{\mathrm{sing}}$ est rare en $Y$.
\mspage{18}
